% Shared module: structured fit (grouping the modes into factors).
\subsection{Grouping the modes: the structured fit}\label{sec:alg-structfit}
The natural route, fitting one polynomial $P$ through all points $(\lambda_p,i\kappa_p)$ and then factoring
it, is ill-posed: a numerical polynomial with noisy coefficients is almost never reducible, and when it is
the factors move discontinuously with the coefficients. In our tests this route split the irreducible
Klein--Gordon and telegraph symbols into spurious factors at noise $10^{-4}$.

We group and fit at the same time. For a structure $\Mfit$ with $N$ factors we solve the
\emph{$N$-varieties} problem, a version of $k$-means clustering in which the cluster centres are curves instead of
points: assign every mode to the factor whose variety is nearest, refit each factor to the modes assigned to it, and
repeat the two steps until the assignment no longer changes (8 random starts, 30 iterations). The distance of a mode to factor $Q_j$ is
the first-order distance to its variety, scaled by the rate uncertainty,
\begin{equation}\label{eq:vardist}
  r_j(p)=\frac{\abs{Q_j(\lambda_p,i\kappa_p)}}{\sigma_{\lambda,p}\,\abs{\partial_\tau Q_j(\lambda_p,i\kappa_p)}},
\end{equation}
which plays the role of the Euclidean distance in $k$-means, and the misfit of a mode is $\min(r_j(p)^2,9)$; the cap makes isolated outliers count as a fixed
penalty instead of dragging a factor away. This is a clustering problem in which the clusters are
algebraic curves rather than points: the grouping of \cref{thm:factor} learned directly from the
spectrum.
